\section{Related Work}
\label{sec:related}

Our work sits at the intersection of three research streams: reward misspecification and overoptimization in RLHF, annotation quality and variance, and automation bias in human judgment. Prior work in each stream addresses downstream symptoms of annotation drift --- benchmark degradation, label noise, evaluation criterion shift --- but not the upstream source in the annotation pipeline itself.

\subsection{Reward Misspecification and RLHF Overoptimization}

Reinforcement Learning from Human Feedback~\citep{Christiano2017Deep,Ouyang2022Training} trains reward models on human preference labels and uses these to fine-tune language models via policy gradient methods. A central challenge is reward misspecification: reward models learn to optimize a proxy for quality rather than quality itself, leading to Goodhart's Law violations as optimization pressure increases~\citep{Gao2022Scaling}.

\citet{Pan2022Effects} provide a systematic framework for mapping reward misspecification to benchmark degradation, demonstrating that even small misspecification errors compound over RLHF training rounds to produce measurable accuracy drops on TruthfulQA and BIG-Bench tasks. \citet{Coste2023Reward} show that annotation variance --- disagreement between annotators --- limits reward model reliability and contributes to overoptimization. Both lines of work treat annotator preferences as stationary noise rather than as potentially directionally drifting signals.

\textbf{Gap:} Neither \citet{Pan2022Effects} nor \citet{Coste2023Reward} ask whether the annotation signal itself shifts direction over time. Our work fills this gap by analyzing preference signal stationarity directly.

\subsection{RLHF Annotation Dynamics and Dataset Structure}

The datasets underlying our experiments have been documented by their creators. \citet{Bai2022Training} describe the Anthropic HH-RLHF dataset as a product of sequential annotation phases --- helpful, harmless, and red-team phases --- with annotators evaluating pairs of AI-generated responses. \citet{Stiennon2020Learning} describe the OpenAI WebGPT comparisons dataset, collected via a multi-session crowdwork design. Neither paper analyzes whether annotator preferences shift directionally across phases.

\citet{Ziegler2019Fine} study the effect of reward model scale on RLHF quality but assume preference stationarity throughout. \citet{Perez2022Red} analyze annotation quality in red-teaming datasets but focus on adversarial robustness rather than temporal drift.

\textbf{Gap:} Existing work takes multi-round annotation datasets as static corpora. Treating rounds as temporal strata and asking whether the preference signal changes direction across strata is a methodological contribution of this paper.

\subsection{LLM-as-Judge and Evaluation Adaptation}

The closest work to ours is \citet{Thakur2024Judging}, who document criterion shift in LLM-as-judge evaluation: language model evaluators adapt their assessment criteria when exposed to AI-generated text, becoming more favorable toward AI-typical stylistic features over time. Our work extends this to human annotators in RLHF \emph{training} pipelines.

\citet{Liang2023Holistic} and related work on verbosity bias in LLM evaluation document that AI evaluators prefer longer, more structured responses independent of quality. We show analogous preferences emerging in human annotation patterns across rounds.

\subsection{Automation Bias and Human-AI Interaction}

Automation bias~\citep{Skitka1999Accountability,Parasuraman2010Complacency} --- the tendency to favor suggestions from automated systems --- is one of the most replicated findings in human factors research. The effect is strongest under conditions of decision uncertainty, time pressure, and high cognitive load.

\citet{Lee2004Trust} provide a comprehensive review of trust in automation, documenting criterion drift even in high-stakes professional settings. \citet{Dietvorst2015Algorithm} show that algorithmic aversion can reverse to appreciation under repeated exposure, suggesting that initial human skepticism of AI style may erode across annotation rounds.

\subsection{Positioning of This Work}

Our work occupies a previously empty cell in the $2\times2$ defined by \{training data vs.\ evaluation\} $\times$ \{measuring drift vs.\ correcting drift\}. The AAI provides the first scalar, pre-registerable instrument for measuring drift in training data annotation.
